\chapter{METHODOLOGY}

This chapter details the systematic approach undertaken to investigate Partial Discharge phenomena. It encompasses the overall research design, the acquisition and nature of datasets, the methodologies for feature extraction from partial discharge signals, the classification techniques employed for identifying partial discharge sources or types, and the metrics used for evaluating the performance of these techniques. A rigorous methodology is paramount for ensuring the reliability and validity of research findings in the complex domain of partial discharge analysis.

\section{Overview of the Proposed Methodology}

The methodology for partial discharge research generally follows a structured pipeline. This typically begins with the acquisition or sourcing of partial discharge signal data, often from experimental setups or simulations mimicking real-world conditions such as defects in power cables, transformers, or \parencite{Borghei_2022_1}. These signals are frequently contaminated with noise, necessitating robust signal processing and denoising techniques \parencite{Friebe_Jenau_2024, Ghosh_2020, Sun_2019_1}. Subsequently, relevant features are extracted from the cleaned signals in various domains (time, frequency, time-frequency) to create a discriminative representation of the discharge event \parencite{Bag_2019_1}. These features then serve as input to machine learning classifiers or other analytical models to identify partial discharge types, locate their sources, or assess insulation degradation \parencite{Aldosari_2023, Anjum_2017_1, Rauscher_2024}. Finally, the performance of these models is rigorously evaluated using appropriate metrics to ascertain their accuracy and reliability in practical diagnostic scenarios \parencite{Srivastava_2023}. Building upon this established framework, the present study implements a specific methodology designed for the comparative analysis of traditional machine learning algorithms. This research leverages a unique, long-term dataset acquired through the contactless antenna-based monitoring of partial discharges on covered overhead power lines \parencite{Klein_2023_1, Klein_2024}. A core component of our approach is the extraction of a comprehensive suite of features from multiple domains, a strategy widely advocated for enhancing diagnostic accuracy \parencite{Lu_2020_2}. These features are then systematically used to train and test a selection of established classifiers in a comparative framework \parencite{Amaya_Sanchez_2024}. The ultimate goal is to identify optimal feature-algorithm pairings for this specific application and to critically assess the generalizability of the resulting models, thereby providing insights into robust diagnostic solutions for real-world field conditions \parencite{Seitz_2024}. The subsequent sections of this chapter will detail each of these stages: the dataset and preprocessing steps (Section 3.2), the feature engineering process (Section 3.3), the selection and implementation of machine learning classifiers (Section 3.4), and the protocol for performance evaluation (Section 3.5).

\section{Datasets}

The research on Partial Discharge (PD) analysis and detection is fundamentally
driven by data, and this work specifically utilizes a unique dataset from the long-term,
on-line monitoring of medium-voltage (MV) overhead power lines. This approach aligns
with a growing trend in the field that moves beyond purely laboratory-based exper￾iments to include data captured from in-situ operational equipment, which presents
both distinct challenges and opportunities for developing robust diagnostic systems Kim
et al., 2020; Sikorski and Gielniak, 2025.
The primary dataset for this research was obtained using a contactless antenna
method designed specifically for detecting partial discharges in XLPE-covered conduc-
4
tors. The use of antennas for PD detection is a well-established non-invasive technique
that captures the electromagnetic waves radiated by discharge events, allowing for re￾mote monitoring without physical contact with the high-voltage components Azam et
al., 2024; Cruz et al., 2023. This radiometric approach is particularly advantageous for
monitoring widespread assets like overhead lines, as it circumvents the need for direct
sensor installation on every component. Our focus on overhead lines with covered con￾ductors is supported by a significant body of research that has specifically investigated
fault detection and PD phenomena in these systems, acknowledging their importance
for overall grid reliability and the unique challenges they present K. Chen et al., 2021;
C. Zhang et al., 2024.
The dataset is characterized by its extensive temporal scope and geographical dis￾tribution. It comprises nearly three years of continuous data, collected at one-hour
intervals from nine different measuring stations situated in Czechia and Slovakia. Each
sample within this dataset is a time-series signal representing a 20-millisecond snapshot
of the electromagnetic activity captured by the antenna. This long-term, continuous
data acquisition from multiple sites provides a rich source of information, capturing a
wide variety of PD signals under diverse environmental and operational conditions. The
ability to track pattern evolution over extended periods is crucial for prognostic health
management and anomaly detection Florkowski, 2024; Orellana et al., 2023.
However, a major challenge inherent in such field-measured data is the high level
of ambient noise and interference that can obscure the faint PD signals. The mea￾surements are often contaminated with various disturbances, including Gaussian white
noise, periodic narrowband interference from communication systems, and stochastic
pulse-shaped interference from other electrical sources Bao et al., 2024; Cao et al., 2025.
Consequently, robust signal processing, filtering, and denoising techniques are critical
preliminary steps for any reliable PD detection pipeline, with many studies proposing
methods based on wavelet transforms, singular value decomposition, or other adaptive
filtering approaches to suppress interference Sahnoune et al., 2024; Xu et al., 2023.
A key strength of our dataset is the availability of reliable ground truth informa-
5
tion, which is essential for the development and validation of diagnostic models. The
contactless antenna systems were deployed at the same monitoring stations where a
traditional galvanic contact method, used by power distribution companies, is also in
operation. This parallel measurement provides a direct, industry-accepted reference
for confirming PD activity. Furthermore, the dataset has been enhanced with manually
curated labels provided by human experts. This expert validation is invaluable for training and assessing the performance of supervised machine learning algorithms, ensuring
that the models learn to identify genuine PD events accurately. The availability of such
high-quality, labeled data enables the application of various supervised machine learning
models, such as Support Vector Machines (SVM) and ensemble methods like Random
Forests, to achieve high-accuracy fault detection and classification Govindarajan et al.,
2021; Serttaş and Hocaoǧlu, 2020. The large volume of data also necessitates research
into efficient data handling, such as novel lossy compression methods and deployment
on edge computing devices for real-time, on-site analysis Klein et al., 2024

\section{Feature Extraction}
To facilitate the classification of partial discharge (PD) signals, a critical stage in the
diagnostic pipeline is the transformation of raw, preprocessed time-series data into a
concise and discriminative set of numerical features Kumar et al., 2024. For this study,
a multi-domain feature extraction strategy was employed, encompassing descriptors
from the time, statistical, frequency, and time-frequency domains. This comprehensive
approach aims to construct a robust representation of PD events by capturing the distinct
signal characteristics essential for accurate diagnosis by machine learning classifiers Lu
et al., 2020; Mas’ud et al., 2016. The subsequent sections provide a detailed account
of the features calculated for this purpose.

\subsection{Time-Domain Features}
Time-domain features are derived directly from the individual PD pulse waveforms to
quantify their fundamental physical properties. These features characterize the shape,
magnitude, and temporal occurrence of the pulse, which are intrinsically linked to the
underlying discharge physics and provide crucial information for classification Kumar et
al., 2024; Lu et al., 2020. The analysis of these characteristics enables the differentiation
of various PD phenomena based on how the discharge manifests as a transient electrical
signal Boczar et al., 2017.
The peak amplitude, signifying the maximum magnitude of the pulse, serves as a
primary indicator of the discharge intensity and is directly related to the energy released
during the event K. Zhang et al., 2024. Different PD sources, such as corona, surface
discharges, or internal voids, often generate pulses with distinct amplitude ranges. This
makes amplitude a valuable feature for assessing a defect’s severity and for distinguish￾ing between source types, as the physical mechanism and geometry of the defect govern
the amount of charge mobilized in the discharge Boczar et al., 2017; Martínez-Tarifa et
al., 2022.
The pulse rise time, defined as the interval for the signal to increase from 10%
to 90% of its peak value, reflects the speed of the ionization process Lu et al., 2020.
PD pulses are characterized by very rapid rise times, often in the nanosecond or sub￾nanosecond range Darwish et al., 2019. This parameter is highly effective for dis￾criminating between PD sources because phenomena such as streamer-like corona dis￾charges exhibit different development speeds compared to Townsend-like discharges in
internal voids Xavier et al., 2021. Consequently, rise time is a critical feature, though
it can be affected by signal degradation and dispersion during propagation from the
source to the sensor Lu et al., 2020.
Pulse duration, frequently measured as the full width at half maximum (FWHM),
quantifies the time for which the discharge event persists and is related to the recom￾bination time of charge carriers following the initial ionization Martínez-Tarifa et al.,
2022. This temporal characteristic is another key descriptor of the pulse shape that aids
7
in distinguishing between defect types. Specific discharge mechanisms, influenced by
factors like gas pressure and electrode geometry, produce characteristically longer or
shorter pulses, making duration a useful parameter for classification Boczar et al., 2017;
Uwiringiyimana et al., 2022.
Finally, the pulse count per cycle tallies the total number of PD events detected
within one period of the power frequency, directly quantifying the repetition rate and
activity level of the PD source Gillis et al., 2023. A high pulse count can indicate a more
active or severe insulation defect, and monitoring this metric is essential for tracking the
evolution and stability of a fault over time Ma et al., 2021. The number of discharges
per cycle, along with their phase distribution, is fundamental to forming Phase-Resolved
Partial Discharge (PRPD) patterns, which are widely used for identifying specific defect
types Boczar et al., 2017; Kasinathan et al., 2017

\subsection{Statistical Features}
Statistical features provide a quantitative description of the shape and distribution
of discharge activity within the Phase-Resolved Partial Discharge (PRPD) pattern. These
descriptors are derived from the distribution of pulse magnitudes or counts across the
phase angles of the AC voltage cycle. For this work, the higher-order statistical moments
of skewness and kurtosis were computed, as they are particularly effective for captur￾ing the nuanced characteristics of PRPD patterns that differentiate various discharge
phenomena Boczar et al., 2017.
Skewness measures the degree of asymmetry in the distribution of PD pulses rel￾ative to the center of the phase axis. This feature is valuable because the physical
mechanisms governing PD can differ significantly between the positive and negative
half-cycles of the applied voltage, leading to characteristically asymmetric PRPD pat￾terns Boczar et al., 2017. For instance, corona from a sharp point electrode typically
produces a much higher discharge activity during one polarity of the AC cycle, resulting
in a highly skewed distribution Yao et al., 2016. By quantifying this asymmetry, skew-
8
ness provides a powerful metric for distinguishing between sources like corona, surface,
and internal void discharges, each of which exhibits unique symmetry properties in its
PRPD profile Bohari et al., 2021.
Kurtosis quantifies the ”peakedness” or ”tailedness” of the distribution, indicating
whether the PD pulses are concentrated within narrow phase windows or are more
broadly dispersed across the cycle. A high kurtosis value signifies a distribution with
sharp, narrow peaks and heavy tails, which is characteristic of discharge types that occur only at specific, repeatable moments near the voltage peaks Boczar et al., 2017.
Conversely, a lower kurtosis suggests a flatter, more uniform distribution of pulses.
Since different insulation defects generate discharges with distinct phase concentrations, kurtosis serves as a key feature for identifying the nature of the PD source by
characterizing the shape of its statistical distribution Kunicki et al., 2018. Together,
these statistical moments transform the qualitative visual information of a PRPD pattern into a compact and discriminative numerical feature vector suitable for automated
classification Bohari et al., 2021

\subsection{Frequency-Domain Features}
requency-domain analysis is employed to investigate the spectral composition of
partial discharge signals, which provides critical diagnostic information not apparent in
the time domain. By applying a Fast Fourier Transform (FFT) to each captured wave￾form, the signal is converted from its temporal representation to its constituent frequen￾cies X. Li et al., 2018. The significance of this analysis lies in the fact that PD pulses
are inherently fast, transient events with rapid rise times, meaning they are broadband
phenomena containing significant energy across a wide spectrum that can extend into
the High Frequency (HF), Very High Frequency (VHF), and Ultra-High Frequency (UHF)
ranges Darwish et al., 2019.
The frequency signature of a measured PD signal is shaped by two key factors:
the physical nature of the discharge source and the propagation path the electromag-
9
netic wave travels to the sensor Azirani et al., 2021. The equipment itself, such as
a power transformer or a Gas-Insulated Switchgear (GIS) enclosure, acts as a complex waveguide, causing frequency-dependent attenuation and reflection that alters the
signal’s spectral characteristics Promwong and Tiengthong, 2020. Consequently, the
energy distribution across different frequency bands becomes a unique fingerprint for a
specific defect type within a given apparatus.
Different types of insulation defects are known to manifest distinctively in the frequency domain. The UHF band (300 MHz–3 GHz) is particularly valuable for PD detection in transformers and GIS, as signals in this range can propagate efficiently within
the metallic enclosures while being less susceptible to external corona and broadcast
interference Salah et al., 2022. For instance, the spectral content of UHF signals can be
analyzed to diagnose and distinguish between specific fault types in GIS, such as those
caused by free-moving particles or fixed protrusions on conductors Ren et al., 2021.
Selecting an optimal frequency band for analysis is therefore crucial for maximizing the
signal-to-noise ratio and detection sensitivity Azirani et al., 2021. For this study, the
spectral power was computed within the distinct HF, VHF, and UHF bands to create a
feature that captures this discriminative energy distribution, providing a robust basis for
classification Maneerot et al., 2019.

\subsection{Time-Frequency Features}
Partial discharge signals are fundamentally transient and non-stationary phenom￾ena, meaning their frequency content evolves rapidly over their short duration. This
characteristic makes time-frequency analysis methods, which can capture both tempo￾ral and spectral information simultaneously, particularly well-suited for their characteri￾zation X. Li et al., 2018. Among these methods, the Discrete Wavelet Transform (DWT)
is exceptionally effective and widely employed in PD analysis for its dual capability of
signal denoising and feature extraction Lu et al., 2020. The DWT decomposes a sig￾nal into various frequency sub-bands at different resolutions, allowing for the effective
10
isolation of PD-related components from background noise and interference Sahnoune
et al., 2024. The efficacy of the DWT is highly dependent on the choice of the mother
wavelet, as its shape should ideally correlate with the morphology of the PD pulses being analyzed to achieve optimal decomposition Sahnoune et al., 2024. For this study,
the Daubechies 4 (db4) wavelet was selected. The Daubechies family of wavelets is
frequently chosen for its properties of compact support and orthogonality, and the db4
wavelet, in particular, has been demonstrated in the literature to be effective for the
denoising and analysis of PD signals, thereby representing an established and validated
choice Wang et al., 2020. The signal was decomposed into four levels, which separates
the signal into a series of detail coefficients representing successively narrower highfrequency bands, and a final approximation coefficient representing the low-frequency
content. For feature extraction, the energy of the detail coefficients at each of the four
decomposition levels was calculated. The resulting four-dimensional vector of wavelet
energies serves as a powerful feature set. This approach is justified because the energy
of the coefficients at a specific level quantifies the intensity of the signal within that
particular frequency band Yang et al., 2024. Since different types of PD sources (e.g.,
surface, corona, void) generate pulses with distinct spectral signatures, their energy is
distributed differently across these frequency bands Gao et al., 2023. Therefore, the
vector of wavelet energies creates a compact and highly discriminative signature that
effectively characterizes the unique time-frequency energy distribution of a PD event,
making it an excellent input for machine learning classifiers Yang et al., 2024.

\section{Classification and Cross-Validation}

This section outlines the classification algorithms employed in this study and the strategies used to validate their performance. Emphasis is placed on selecting methods that align with the complex and diverse nature of partial discharge data, followed by a cross-validation framework to ensure rigorous model assessment.

\subsection{Algorithm Selection and Justification}

This study implemented and compared a curated set of five supervised machine
learning algorithms to diagnose insulation defects based on the extracted partial dis￾charge features. The selected models represent a diverse range of classification strate￾gies, from kernel-based methods and tree-based ensembles to instance-based and neu￾ral network approaches, ensuring a comprehensive comparative analysis.
This study implemented and compared a curated set of five supervised machine
learning algorithms to diagnose insulation defects based on the extracted partial dis￾charge features. The selected models represent a diverse range of classification strate￾gies, from kernel-based methods and tree-based ensembles to instance-based and neu￾ral network approaches, ensuring a comprehensive comparative analysis.
The Support Vector Machine (SVM) was included for its capacity to construct op￾timal decision boundaries in high-dimensional feature spaces. By mapping features to
a higher dimension and identifying a separating hyperplane that maximizes the mar￾gin between classes, the SVM can effectively model complex decision boundaries. Its
application in PD diagnostics spans numerous detection modalities; it has been used
to identify ultra-high-frequency (UHF) signals in Gas-Insulated Switchgear (GIS) from
moment-based features Bin et al., 2020 and to localize PD sources when integrated with
image processing workflows X. Li et al., 2019.
This adaptability makes the SVM a valuable tool for PD analysis across different
insulation systems and data types. Research has documented its use for classifying dis￾charges in traditional oil-paper insulation systems Woon et al., 2016 and for interpreting
abstract feature sets, such as those derived from the fractal analysis of PD images Basha￾ran et al., 2018. A comparative study by Dobrzycki et al. (2019) further underscores
its utility by evaluating it alongside ANNs for the detection of acoustic emission sig￾nals accompanying electrical treeing, confirming its relevance in multi-modal diagnostic
scenarios Dobrzycki et al., 2019.
As a representative of ensemble learning, the Random Forest classifier was imple￾mented. This algorithm constructs a multitude of decision trees during training and
12
outputs the mode of their individual predictions, a process that inherently mitigates the
risk of overfitting to the noisy and variable measurements typical of PD phenomena. Its
direct application in the acoustic identification of PD sources in power equipment illus￾trates its foundational value in this field Zhu et al., 2025. The model’s integration with
advanced signal processing, such as using the S-Transform to analyze optical sensor
data for PD location, points to its role within more sophisticated diagnostic frameworks
Bag et al., 2019. The development of specialized variants, like hypergraph-based Ran￾dom Forest algorithms, further highlights its established position in improving PD pattern
classification Govindarajan et al., 2021.
The K-Nearest Neighbors (KNN) algorithm offers a non-parametric, instance-based
learning approach. Its core principle—classifying a data point based on the majority
class of its nearest neighbors—avoids making prior assumptions about the underlying
data distribution. This is a key justification for its inclusion, as the statistical properties
of PD signals from different insulation defects are often complex and unknown Lu et
al., 2020. Its application has been demonstrated for the robust classification of partial
discharges in various insulation systems Woon et al., 2016. Beyond direct classification,
its logic has been integrated into more intricate diagnostic systems, such as matching
new PD data within a case-based reasoning framework for GIS fault analysis Dai et al.,
2019.
To establish an interpretable benchmark, a single Decision Tree classifier was im￾plemented. The primary value of this model in an engineering diagnostic context is
its transparent, hierarchical structure, which renders the classification logic as a set of
explicit, human-readable rules Lee and Horng, 2020. This interpretability is critical for
validating diagnostic outcomes and justifying maintenance actions. While a single tree
can be less robust than ensemble methods, it serves two vital functions: it provides a
clear performance baseline, as seen in transformer fault diagnosis studies Maseko et al.,
2025, and it constitutes the fundamental building block for the more complex Random
Forest classifier Prabhu et al., 2016.
Finally, an Artificial Neural Network (ANN) was included to model the complex, non-
13
linear relationships inherent in PD phenomena. As computational models capable of
learning intricate patterns from data, ANNs are essential for capturing the link between
a measured signal and its underlying physical defect. A comprehensive survey by Mas’ud
et al. (2016) confirms their well-established role in PD recognition Mas’ud et al., 2016.
The breadth of their applicability is captured in studies that assess their performance in
various contexts, from diagnosing PD in motor stator windings Y.-T. Chen et al., 2018 to
identifying defect patterns in cables under DC stress Khan and Koo, 2017.

\subsection{Model Validation Strategy}

To ensure a robust and unbiased evaluation of model performance, a k-fold crossvalidation strategy was employed. This validation technique is essential for mitigating the risk of overfitting, particularly with limited experimental datasets, and for providing a reliable estimate of each model's generalization capability on unseen data. It is considered a standard procedure for assessing machine learning models in PD monitoring systems \parencite{Srivastava_2023} and a critical component in the qualification process for PD analysers \parencite{Vera_2023}. In this study, 10-fold cross-validation was utilized. The dataset was partitioned into ten mutually exclusive subsets, or 'folds', of approximately equal size. The classification process was then iterated ten times; in each iteration, one distinct fold was reserved as the test set for validation, while the remaining nine folds were used for training the model. The final performance metrics for each algorithm were calculated by averaging the results obtained across all ten folds. This approach ensures that the reported performance is a stable and comprehensive reflection of the model's diagnostic capabilities.

\section{Evaluation Metrics}

A comprehensive and rigorous evaluation of classifier performance is essential for developing reliable diagnostic systems. The assessment of the proposed partial discharge (PD) classification models will be conducted using a suite of standard metrics, each chosen to provide a distinct perspective on the model's effectiveness. This multifaceted approach ensures that the evaluation captures not only the overall correctness but also the specific strengths and weaknesses related to identifying individual fault types, as well as the practical feasibility for deployment in monitoring systems.

\subsection{Confusion Matrix}

The confusion matrix is a fundamental tool that provides a detailed breakdown of a classifier's performance on a class-by-class basis. It is structured as a table that explicitly visualizes the counts of correct and incorrect predictions, tabulating the number of true positives (TP), true negatives (TN), false positives (FP), and false negatives (FN) for each diagnostic category. This matrix serves as the bedrock for calculating most other performance metrics and offers a direct, qualitative analysis of misclassifications \parencite{Sikorski_2025}. In the complex context of PD diagnosis, the confusion matrix is invaluable because it reveals not just that errors occurred, but specifically which defect types are being confused with one another. This granular insight is critical for understanding the diagnostic limitations of a model and for guiding future improvements, such as refining feature extraction to better distinguish between similar fault signatures.

\subsection{Accuracy}

Accuracy represents the proportion of all correctly classified instances relative to the total number of instances evaluated, calculated as (TP + TN) / (TP + TN + FP + FN). It serves as a primary, high-level indicator of a model's overall performance and is a fundamental benchmark metric used extensively throughout PD classification research \parencite{Aldosari_2023}. While it provides an intuitive and useful general summary of correctness, its diagnostic utility can be limited, particularly when dealing with the imbalanced class distributions that are characteristic of fault diagnosis data. In such scenarios, a model can achieve a high accuracy score by simply predicting the majority class (e.g., 'no fault'), while failing to identify rare but critical fault conditions. Therefore, it is always interpreted in conjunction with more nuanced, class-specific metrics.

\subsection{Precision}

Calculated as TP / (TP + FP), precision quantifies the exactness of the classifier by measuring the proportion of instances classified as positive that are genuinely positive. For PD diagnostics, high precision is essential for ensuring the reliability and trustworthiness of a fault warning. This metric directly correlates to minimizing the false alarm rate, which is a critical operational objective. \subsection{Recall (Sensitivity)}

Calculated as TP / (TP + FN), recall, also known as sensitivity or the true positive rate, measures the classifier's ability to correctly identify all actual instances of a specific fault class. Within the domain of high-voltage asset management, achieving a high recall is paramount, as it reflects the model's completeness in fault detection. Failing to detect a genuine PD source---a false negative---can have severe consequences, potentially leading to progressive insulation degradation, unexpected outages, and ultimately, catastrophic equipment failure \parencite{Adhikari_2024}. This metric is therefore a critical indicator of the system's effectiveness in proactive risk mitigation and ensuring the safety and reliability of the power network.

\subsection{F1-Score}

The F1-Score is computed as the harmonic mean of precision and recall, using the formula 2 * (Precision * Recall) / (Precision + Recall). It provides a single, balanced measure of a model's performance by considering both false positives and false negatives. It is particularly valuable for evaluating classifiers on imbalanced datasets, which is a common scenario in fault diagnosis where data from certain critical fault types may be scarce compared to normal operating data \parencite{Zemouri_2023}. 